\section{O que foi feito}

Com o crescimento da inteligencia artificial, a necessidade de processamente de dados massivos vem se tornando cada vez mais relevante. 
Em conjuntos de dados de grande escala eh comum ser necessario representar o relacionamente entre diferentes entidades. 
É natural que esse tipo de representacao seja feito atraves de modelagens em grafo.

Em uma rede social, vertices podem representar usuarios e arestas podem representar
amizades, seguidores, mensagens ou interacoes \cite{easleykleinberg}. Na Web,
paginas sao vertices e hiperlinks sao arestas \cite{brinpage}. Em seguranca
computacional, grafos podem modelar fluxos de rede, dependencias entre
processos, tentativas de acesso e relacoes entre entidades suspeitas
\cite{akoglugraphanomaly}. Em biologia, vertices podem representar proteinas,
genes ou moléculas, e arestas podem representar interacoes
\cite{barabasinetworkbiology}. Em todos esses casos, a Busca em Largura aparece como uma
primitiva para calcular alcancabilidade, distancia em numero de arestas e
expansao de vizinhanca.


Este trabalho propoem um novo metodo para o algoritmo de busca em largura no
contexto de dados massivos explorando metodos de paralelismo, com o foco em processadores manycore.  A linguagem CUDA em conjunto com a biblioteca de comunicacao da NVIDIA NVSHMEM,
foram utilizados para o desemvolvimento do algoritimo, os testes foram feitos em dois ambientes:
\begin{itemize}
\item Sistema operacional Linux, GPU NVIDIA 5060ti 16GB, e processador Ryzen 5 5600x, 32GB de memoria RAM.
\item No de uma GB200, linkar. \leal{ explicar melhor esse ambiente}
\end{itemize}

% O trabalho utiliza o Graph500 como benchmark para avaliacao da performance.
% 
% A literatura de sistemas de grafos trata a BFS como uma das operacoes centrais
% para avaliar plataformas de processamento de dados irregulares. O Graph500 foi
% proposto para complementar benchmarks numericos tradicionais com uma carga de
% trabalho orientada a dados, motivado pela diferenca entre simulacoes fisicas,
% que tendem a ser estruturadas e intensivas em ponto flutuante, e aplicacoes de
% informatica, que tendem a ser nao estruturadas, orientadas a inteiros e com
% baixa localidade \cite{murphygraph500}. Sua especificacao define a BFS como o
% Kernel 2 do benchmark, executado sobre grafos sinteticos de escala crescente, e
% usa a metrica TEPS, ou arestas atravessadas por segundo, para medir desempenho
% \cite{graph500spec}. Isso mostra que a BFS nao e apenas um exercicio teorico:
% ela e usada como proxy para medir a capacidade de supercomputadores processarem
% analise de grafos massivos.
% 
% Em redes sociais, a BFS modela consultas de vizinhanca e distancia social. A
% pergunta ``quais usuarios estao a ate dois ou tres saltos de uma conta?'' e uma
% execucao limitada de BFS. Esse tipo de consulta aparece em recomendacao de
% amizades, propagacao de informacao, deteccao de comunidades locais e analise de
% influencia \cite{easleykleinberg}. Em grafos com bilhoes de relacoes, a
% expansao de poucos niveis pode atingir uma quantidade enorme de vertices,
% especialmente quando existem hubs de grau muito alto. Por isso, mesmo uma BFS
% limitada por profundidade pode exigir processamento paralelo e estruturas de
% dados distribuidas.
% 
% Em grafos da Web, a BFS tambem e uma abstracao natural. Um rastreador
% \textit{web crawler} pode iniciar em um conjunto de paginas sementes e expandir
% links por camadas, descobrindo novas paginas a cada nivel \cite{brinpage}. O
% sistema Pregel, do Google, foi proposto para processamento de grafos em larga
% escala e apresenta algoritmos como caminhos minimos, PageRank, propagacao de
% rotulos e outras computacoes iterativas sobre grafos distribuidos \cite{pregel}.
% Embora Pregel seja um modelo mais geral que BFS, sua execucao em superpassos e
% troca de mensagens entre vertices e muito proxima da ideia de fronteiras
% sucessivas usada em BFS distribuida.
% 
% Em seguranca, investigacao de fraudes e analise de dependencias, BFS e usada
% para expandir a vizinhanca de uma entidade de interesse. A partir de uma conta,
% um endereco IP, uma maquina infectada ou um processo suspeito, a busca em
% largura identifica entidades conectadas por poucos saltos. Essa forma de
% exploracao e importante porque muitas relacoes relevantes nao aparecem em uma
% aresta direta, mas em cadeias curtas: contas que compartilham dispositivos,
% maquinas que se comunicam com o mesmo servidor, ou processos que dependem dos
% mesmos arquivos \cite{akoglugraphanomaly}. Em grafos corporativos ou de
% telecomunicacoes, essas buscas podem envolver centenas de milhoes de vertices e
% precisam responder em tempo interativo ou quase interativo.
% 
% O desafio e que os grafos modernos podem ser massivos e irregulares. Grafos de
% redes sociais e da Web frequentemente possuem distribuicao de grau altamente
% desbalanceada: poucos vertices concentram muitas arestas, enquanto a maioria tem
% grau baixo. Essa propriedade gera duas dificuldades para BFS. Primeiro, ha
% desbalanceamento de trabalho, pois expandir um vertice de alto grau custa muito
% mais do que expandir um vertice comum. Segundo, ha muita contencao, pois muitos
% caminhos podem tentar descobrir o mesmo vertice ao mesmo tempo.
% 
% Alem disso, a BFS tem baixa localidade de memoria. Em matrizes densas ou
% vetores, os acessos costumam ser contiguos e previsiveis. Em grafos, ao
% expandir a fronteira, o algoritmo segue listas de adjacencia que podem apontar
% para posicoes distantes e aparentemente aleatorias. Consequentemente, a BFS e
% normalmente limitada por largura de banda, latencia de memoria, atomicos,
% sincronizacao e comunicacao entre maquinas, e nao por operacoes aritmeticas.
% 
% Trabalhos em arquiteturas paralelas mostram como essas dificuldades aparecem na
% pratica. Beamer, Asanovic e Patterson propuseram a BFS \textit{direction
% optimizing}, que alterna entre uma expansao top-down, baseada na fronteira
% atual, e uma expansao bottom-up, que examina vertices ainda nao visitados
% procurando vizinhos na fronteira \cite{beamerbfs}. Essa mudanca reduz trabalho
% em grafos de pequeno diametro e grande fronteira, um caso comum em grafos
% sociais e da Web. Merrill, Garland e Grimshaw estudaram BFS em GPU e mostraram
% que a primitiva e representativa de computacoes paralelas com acesso a memoria e
% distribuicao de trabalho irregulares; sua implementacao alcançou bilhoes de
% arestas atravessadas por segundo em uma ou mais GPUs \cite{merrillgpu}.
% 
% Portanto, BFS aparece em problemas reais de duas formas complementares. Como
% algoritmo final, ela responde consultas de alcance, distancia e vizinhanca. Como
% primitiva de sistema, ela estressa exatamente os componentes que tornam grafos
% massivos dificeis: memoria irregular, fronteiras dinamicas, atomicos,
% desbalanceamento e comunicacao. Por isso, estudar BFS em GPU e em ambientes
% distribuidos e uma forma direta de estudar o desempenho de analise de grafos em
% escala real.
